%% ============================================================
%% Academic Journal Paper Template
%% Compatible with: IEEE, ACM, Elsevier, Springer conventions
%% ============================================================

\documentclass[10pt, twocolumn]{article}

%% ---- Core Packages ----------------------------------------
\usepackage[T1]{fontenc}
\usepackage[utf8]{inputenc}
\usepackage{lmodern}
\usepackage{microtype}           % Microtypographic refinements

%% ---- Page Geometry ----------------------------------------
\usepackage[
  top=1in, bottom=1in,
  left=0.75in, right=0.75in,
  columnsep=0.25in
]{geometry}

%% ---- Mathematics ------------------------------------------
\usepackage{amsmath, amssymb, amsthm}
\usepackage{mathtools}
\usepackage{bm}                  % Bold math symbols

%% ---- Figures & Tables -------------------------------------
\usepackage{graphicx}
\usepackage{booktabs}            % Professional table rules
\usepackage{multirow}
\usepackage{array}
\usepackage{caption}
\usepackage{subcaption}
\usepackage{float}

%% ---- Algorithms -------------------------------------------
\usepackage[ruled, vlined, linesnumbered]{algorithm2e}

%% ---- Code Listings ----------------------------------------
\usepackage{listings}
\usepackage{xcolor}

\lstdefinestyle{codestyle}{
  backgroundcolor=\color{gray!8},
  basicstyle=\ttfamily\footnotesize,
  breakatwhitespace=false,
  breaklines=true,
  captionpos=b,
  commentstyle=\color{green!50!black},
  keywordstyle=\color{blue}\bfseries,
  stringstyle=\color{orange!70!black},
  numberstyle=\tiny\color{gray},
  numbers=left,
  numbersep=5pt,
  showstringspaces=false,
  frame=single,
  rulecolor=\color{gray!40},
  tabsize=2
}
\lstset{style=codestyle}

%% ---- Hyperlinks -------------------------------------------
\usepackage[
  colorlinks=true,
  linkcolor=blue!70!black,
  citecolor=green!50!black,
  urlcolor=blue!60!black
]{hyperref}
\usepackage{url}

%% ---- Bibliography -----------------------------------------
\usepackage[numbers, sort&compress]{natbib}

%% ---- Miscellaneous ----------------------------------------
\usepackage{lipsum}              % Lorem ipsum filler — remove in production
\usepackage{enumitem}
\usepackage{siunitx}             % SI units: \SI{3.14}{\mega\hertz}
\usepackage{cleveref}            % Smart cross-references: \cref{fig:foo}

%% ---- Theorem Environments ---------------------------------
\theoremstyle{plain}
\newtheorem{theorem}{Theorem}[section]
\newtheorem{lemma}[theorem]{Lemma}
\newtheorem{corollary}[theorem]{Corollary}
\newtheorem{proposition}[theorem]{Proposition}

\theoremstyle{definition}
\newtheorem{definition}[theorem]{Definition}
\newtheorem{example}[theorem]{Example}

\theoremstyle{remark}
\newtheorem{remark}[theorem]{Remark}

%% ---- Custom Commands --------------------------------------
\newcommand{\R}{\mathbb{R}}
\newcommand{\N}{\mathbb{N}}
\newcommand{\Z}{\mathbb{Z}}
\newcommand{\E}{\mathbb{E}}
\newcommand{\Prob}{\mathbb{P}}
\newcommand{\norm}[1]{\left\lVert #1 \right\rVert}
\newcommand{\abs}[1]{\left\lvert #1 \right\rvert}
\newcommand{\inner}[2]{\left\langle #1,\, #2 \right\rangle}
\newcommand{\ie}{\textit{i.e.}\xspace}
\newcommand{\eg}{\textit{e.g.}\xspace}
\newcommand{\etal}{\textit{et al.}\xspace}

%% ---- Caption Formatting -----------------------------------
\captionsetup{
  font=small,
  labelfont=bf,
  format=hang,
  justification=justified
}

%% ---- Header / Footer --------------------------------------
\usepackage{fancyhdr}
\pagestyle{fancy}
\fancyhf{}
\fancyhead[L]{\small\itshape Ne.app, Vol.~2, No.~1, March 2026}
\fancyhead[R]{\small\thepage}
\renewcommand{\headrulewidth}{0.4pt}

%% ============================================================
%%  FRONT MATTER
%% ============================================================

\title{%
  \vspace{-1.5em}%
  \large\textbf{All you need is Integrals for Debugging.}%
}

\author{%
  \textbf{Amlal El Mahrouss}$^{1}$\thanks{Corresponding author. Email: amlal@nekernel.org}
  \\amlal@nekernel.org, amlalelmahrouss@icloud.com\\[0.4em]
  \small $^{1}$Ne.app Journal
}

\date{%
  \small 3 March 2026
}

%% ============================================================
\begin{document}


\twocolumn[{%
  \maketitle
  \vspace{-0.5em}
  \rule{\linewidth}{0.4pt}
  %% ---- Abstract -------------------------------------------
  \begin{center}
  \begin{minipage}{0.92\linewidth}
    \small
    \textbf{Abstract.}\enspace
    This development presents an approach of software debugging[1] in which we'll name 'Integrals of Debuggging' (ID).
    We propose a mathematical of which integrals are used to analyze, modify, and apply, software patches to programs.

    
    \medskip
    \textbf{Keywords:}\enspace
    Computer Science; Mathematics;
  \end{minipage}
  \end{center}
  \vspace{0.8em}
  \rule{\linewidth}{0.4pt}
  \vspace{1em}
}]

\section{Introduction}

Consider an open interval from $n$ to $p$ over a field $\mathbb{K}$ which describes a program and its properties. What are the techniques in which we can analyze and audit its properties and behavior? We will try to impose a framework in this report.

\section{Definition of the LongReturn(x, t) function}
\label{sec:cl-example-integrals}

Let an integral of $\operatorname{LongReturn(x, t)}$[2] be defined where the value $n > 0$ and $p > 0, \quad n < p$:

\begin{equation}
    \operatorname{LongReturn(x, t)} \coloneqq \int_{n}^{p} \operatorname{FastReturn(x, t)} \cdot d \operatorname{t}
\end{equation}
We assume that $x$ and $t$ are always greater or equal than one, $\operatorname{LongReturn(x, t)}$ shall always be greater than $\operatorname{FastReturn(x, t)}$ for all such value in $\mathbb{K}$.
This introduction of the LongReturn, will let us jump to our next section about its applications in Computer Analysis.

\section{Constraints}

However before starting, we shall provide the following conditions from [2]:
\begin{equation}
    \operatorname{LongReturn(x, t)} \geq \operatorname{FastReturn(x, t)}, \quad \operatorname{x} \geq 1
\end{equation}
The variable x shall also be defined in $\mathbb{K}$ throughout the integral. We shall now proceed into the next section.

\section{Definition of the DerRet(x, t) function}

The following is the derivative of $\operatorname{LongReturn(x, t)}$, which equals to:
\begin{equation}
    \operatorname{DerRet(State, Time\_State)} \coloneqq \frac{\delta \operatorname{Travel(State)}}{\delta Time\_State} + Rem
\end{equation}
Such that the following is assumed:
\begin{equation}
    \quad State, \operatorname{Travel(State)} \in \mathbb{K}
\end{equation}
Indeed, the function $\operatorname{Travel(State)}$ is equal to the program state at the variable $\operatorname{State}$ which represents the program's state.

\section{Applications}

One application of the $\operatorname{LongReturn(x, t)}$ is software debugging, more so, the break-point system of a debugger. The following section includes graphs to measure the efficiency of our approach in such process.

\end{document}


